\documentclass[11pt,a4paper]{article}
\usepackage[margin=25mm,headheight=14pt]{geometry}
\usepackage[T1]{fontenc}
\usepackage{lmodern}
\usepackage{amsmath,amssymb,amsthm,mathtools}
\usepackage{booktabs,array,longtable,microtype}
\usepackage{xcolor,hyperref,fancyhdr}
\definecolor{linkblue}{RGB}{23,57,93}
\hypersetup{colorlinks=true,linkcolor=linkblue,urlcolor=linkblue,citecolor=linkblue,
  pdftitle={The C\_t Route toward Binary Non-Normality of Pi},
  pdfsubject={Formalized analytic progress and an explicit residual-tail criterion}}
\urlstyle{tt}
\pagestyle{fancy}
\fancyhf{}
\fancyhead[L]{\small The $C_t$ route: formalized progress}
\fancyhead[R]{\small Research preprint, v1}
\fancyfoot[C]{\small\thepage}
\setlength{\parskip}{4pt}
\setlength{\parindent}{1.2em}
\setlength{\emergencystretch}{2em}
\newtheorem{theorem}{Theorem}[section]
\newtheorem{lemma}[theorem]{Lemma}
\newtheorem{proposition}[theorem]{Proposition}
\newtheorem{corollary}[theorem]{Corollary}
\theoremstyle{definition}
\newtheorem{definition}[theorem]{Definition}
\theoremstyle{remark}
\newtheorem{remark}[theorem]{Remark}
\newcommand{\N}{\mathbb N}
\newcommand{\R}{\mathbb R}
\newcommand{\C}{\mathbb C}
\newcommand{\ind}{\mathbf 1}
\newcommand{\Normal}{\operatorname{Normal}_2}
\newcommand{\fracpart}[1]{\{#1\}}
\newcommand{\Ct}{C_t}
\newcommand{\repo}{\href{https://github.com/randogelist/PI_NOT_NORMAL_BASE_2}{PI\_NOT\_NORMAL\_BASE\_2}}
\newcommand{\code}[1]{{\small\nolinkurl{#1}}}
\title{\vspace{-12mm}\textbf{The \texorpdfstring{$C_t$}{C_t} Route toward\\Binary Non-Normality of \texorpdfstring{$\pi$}{Pi}}\\[6pt]
  \large Formalized analytic progress and an explicit residual-tail criterion}
\author{\normalsize Research note accompanying \repo}
\date{\normalsize Version 1 \quad$\cdot$\quad 25 October 2026}

\begin{document}
\maketitle
\thispagestyle{plain}
\begin{abstract}
We present the mathematical content and proof boundary of a Lean development
investigating binary non-normality of $\pi$ through a signed dyadic coefficient.
An arctangent representative anchored at the actual constant $\pi$ yields a
logarithmic split, a uniquely determined coefficient $C_t=(-1)^t/4$, and an
explicit second-order remainder. For the actual binary digits, a finite
counting argument bounds the combined frequency of two alternating words by
the empirical tail of a precisely defined parity residual. A sufficiently
small eventual tail then implies non-normality. The required tail estimate
for $\pi$ remains open. We explain why coefficient nonvanishing alone does
not supply it, distinguish historical uses of the notation $C_t$, and map
the statements to a GitHub-hosted Lean verification of 71 source modules
across eight isolated projects. This preprint records verified progress and
its conditional endpoint; it does not assert an unconditional
non-normality theorem.
\end{abstract}

\section{Scope and conventions}
\label{sec:scope}

The retained research route begins with an analytic coefficient and seeks
an estimate on the actual digits of $\pi$. These are separate mathematical
objects. We use $C_t$ only for the real-valued coefficient, and reserve
frequency notation for counts of actual words.

Throughout, $\N=\{0,1,2,\ldots\}$. A Boolean digit is identified with the
real number $0$ or $1$ when it occurs in a sum. For a binary stream $b$,
a word $w\in\{0,1\}^{\ell}$, and $M\ge1$, define
\begin{equation}
 f_M^b(w)=\frac1M\#\{0\le n<M:
       (b_n,\ldots,b_{n+\ell-1})=w\}.
 \label{eq:frequency}
\end{equation}
The count is over $M$ starting positions. A word starting before $M$
may extend beyond position $M-1$; it is not required to fit inside
a truncated prefix.

\begin{definition}
A stream $b$ is binary normal if, for every fixed length $\ell$ and
every $w\in\{0,1\}^{\ell}$,
$f_M^b(w)\longrightarrow 2^{-\ell}$ as $M\to\infty$.
For a real number $x$ we use the fractional binary digits
\begin{equation}
 b_n(x)=\ind_{\{\fracpart{2^n x}\ge1/2\}}.
 \label{eq:bits}
\end{equation}
We write $\Normal(x)$ for normality of this stream.
\end{definition}

The formal development implements words by their binary integer codes;
equality of codes of a fixed length is proved to imply equality of
the constituent bits. Conversely, equal bits give equal codes by the
recursive definition. Its predicate \code{BinaryNormal} is the
fixed-word definition above. The main finite estimate, specialized to
$\pi$, is
\begin{equation}
 F_{2N}(r)\le
 2T_{E_\pi}(r/4,N)+\frac{r+1/2}{N},\qquad N>0,
 \label{eq:mainpreview}
\end{equation}
where $F_M(r)$ counts the two alternating words of length $2r+1$,
and $E_\pi$ is defined in Section~\ref{sec:residual}.
The finite inequality is unconditional. Controlling its right-hand side
as $N\to\infty$ is the remaining distribution problem.

\section{The anchored analytic coefficient}
\label{sec:coefficient}

Let $P(z)=4\arctan_{\C}(z)$, with the complex arctangent used by Mathlib.
The development proves
\begin{equation}
 P(1)=\pi,\qquad
 P(z)=4\sum_{n\ge0}\frac{(-1)^n z^{2n+1}}{2n+1}
 \quad (|z|<1).
 \label{eq:anchor}
\end{equation}
The identity at $1$ follows from $\arctan(1)=\pi/4$.
It does not follow by substituting $z=1$ into the displayed
open-disk power-series statement.

For $M>0$ and $s\ge0$, set $q(M,s)=s/(2M+s)$. Along the relevant
imaginary ray,
\[
 \operatorname{Im}P(iq(M,s))=2\log(1+s/M).
\]
Consequently, the ratio of two readouts from the same representative is
\begin{equation}
 Q(M)=\frac{\operatorname{Im}P(iq(M,1))}
              {\operatorname{Im}P(iq(M,2))}
      =\frac{\log(1+1/M)}{\log(1+2/M)}
      =\frac12+\delta(M),
 \label{eq:split}
\end{equation}
where
\begin{equation}
 \delta(M)=
 \frac{\log\!\left(1+\frac1{M(M+2)}\right)}
      {2\log(1+2/M)}.
 \label{eq:delta}
\end{equation}
The denominator is strictly positive. The last equality follows by
taking logarithms in
$(1+1/M)^2/(1+2/M)=1+1/[M(M+2)]$.
This readout is an analytic ratio; no identification with a digit
probability or a word frequency is made.

\begin{lemma}[Rational bounds for the logarithmic defect]
\label{lem:delta}
For every real $M>0$,
\begin{equation}
 \frac{M}{4(M+1)^2}\le\delta(M)\le\frac1{4M}.
 \label{eq:deltabounds}
\end{equation}
For every $N>0$ and $t\ge0$,
\begin{equation}
 0\le\frac14-N\delta(N+t)\le\frac{t+2}{4N}.
 \label{eq:scaledbounds}
\end{equation}
\end{lemma}
\begin{proof}
For $u\ge0$, $u/(1+u)\le\log(1+u)\le u$.
Applied to the numerator and denominator of \eqref{eq:delta}, these
bounds give
\[
 \frac1{(M+1)^2}\le
 \log\!\left(1+\frac1{M(M+2)}\right)\le\frac1{M(M+2)},
 \qquad
 \frac2{M+2}\le\log(1+2/M)\le\frac2M.
\]
Division yields \eqref{eq:deltabounds}. With $M=N+t$, it follows that
\[
 \frac{N(N+t)}{4(N+t+1)^2}
       \le N\delta(N+t)\le\frac{N}{4(N+t)}\le\frac14.
\]
For the error in the lower bound, after multiplication by positive
denominators, the required inequality is equivalent to
\[
 (t+2)(N+t+1)^2
   -N\bigl((N+t+1)^2-N(N+t)\bigr)\ge0.
\]
Its left-hand side equals
$t(t+2)(N+t)+(2t+3)(N+t)+2t+2$, which is nonnegative.
\end{proof}

Write $N_a=2^{10+a}$ and preserve the signed readout
\begin{equation}
 G(a,t)=\frac12+(-1)^t\delta(N_a+t),\qquad
 H(a,t)=N_a\bigl(G(a,t)-\tfrac12\bigr),
 \quad a,t\in\N.
 \label{eq:dyadic}
\end{equation}

\begin{theorem}[The frozen coefficient and its remainder]
\label{thm:coefficient}
For each fixed $t\in\N$, the limit of $H(a,t)$ as $a\to\infty$ exists
uniquely and equals
\begin{equation}
 C_t=\frac{(-1)^t}{4}.
 \label{eq:Ct}
\end{equation}
In particular, $|C_t|=1/4$ and $C_t\ne0$. For all $a,t\in\N$,
\begin{align}
 |H(a,t)-C_t|&\le\frac{t+2}{4N_a}, \label{eq:Cerror}\\
 \left|G(a,t)-\frac12-\frac{C_t}{N_a}\right|
       &\le\frac{t+2}{4N_a^2}. \label{eq:Gremainder}
\end{align}
If $a\le b$, then
\begin{equation}
 |H(b,t)-H(a,t)|\le\frac{t+2}{4N_a}.
 \label{eq:cauchy}
\end{equation}
\end{theorem}
\begin{proof}
Equation \eqref{eq:Cerror} follows from \eqref{eq:scaledbounds}
and $|(-1)^t|=1$. Since $N_a^{-1}\to0$, the squeeze theorem gives
the limit, and uniqueness follows in $\R$.
Dividing \eqref{eq:Cerror} by $N_a$ gives
\eqref{eq:Gremainder}.
For \eqref{eq:cauchy}, both unsigned values
$N_a\delta(N_a+t)$ and $N_b\delta(N_b+t)$ lie in the interval
$[1/4-(t+2)/(4N_a),\,1/4]$. Their difference is bounded by the
length of that interval.
\end{proof}

In Lean, \code{Pi.C_T t} is obtained from proved unique existence of
the limit for the concrete representative; its value is then identified
by \code{PI_CT_VALUE}. The concrete instance supplies the analytic
existence proof. No unresolved analytic premise is passed into that
instance. Every limit here fixes $t$ before letting $a$ grow.

\section{Actual digits and the analytic parity residual}
\label{sec:residual}

The formal stream \code{piBits} is defined from
$\operatorname{Re}P(1)$. The anchor in \eqref{eq:anchor} gives, for every
$n\in\N$,
\[
 \text{\code{piBits}}(n)=b_n(\pi).
\]
This binds the subsequent word counts to the actual real number $\pi$.

For a binary stream $b$ and a real step sequence $a=(a_j)_{j\ge0}$,
define
\begin{equation}
 S_b(m)=\sum_{j<m}(b_{2j}-b_{2j+1}),\qquad
 A(m)=\sum_{j<m}a_j,\qquad E_{b,a}(m)=S_b(m)-A(m).
 \label{eq:residual}
\end{equation}
For the analytic specialization, the step is
\begin{equation}
 a_j=
 \begin{cases}
 0,&j<512,\\
 \delta(2j)+\delta(2j+1),&j\ge512.
 \end{cases}
 \label{eq:analyticstep}
\end{equation}
The cutoff avoids zero denominators and retains the original dyadic
scale convention. Lemma~\ref{lem:delta} implies
$0\le a_j\le1/2$ for every $j$.
We abbreviate $E_{b(\pi),a}$ by $E_\pi$.

The analytic drift is also quantitatively close to its harmonic
approximation. Put $h_j=0$ for $j<512$ and
$h_j=1/(8j)+1/[4(2j+1)]$ otherwise.
The source proves
\begin{equation}
 0\le\sum_{j<m}h_j-A(m)
 \le \frac1{2048}-\frac1{4\max(m,512)}
 \le\frac1{2048}.
 \label{eq:harmonic}
\end{equation}
Indeed, \eqref{eq:deltabounds} gives the scalar estimate
\[
 0\le\frac1{4M}-\delta(M)
 \le\frac12\left(\frac1M-\frac1{M+1}\right).
\]
Adding it at $M=2j$ and $M=2j+1$ bounds the step error by
$1/(4j)-1/[4(j+1)]$, and the sum telescopes.
This controls two analytic sequences. It does not bound the residual
of either sequence against the actual parity sum $S_{b(\pi)}$.

\section{A finite alternating-word inequality}
\label{sec:finite}

For any residual $E$ and $N\ge1$, define its empirical tail by
\begin{equation}
 T_E(h,N)=\frac1N\#\{0\le j<N:|E(j)|\ge h\}.
 \label{eq:tail}
\end{equation}
For $r\in\N$, let $u_r=(01)^r0$ and $v_r=(10)^r1$, and write
\begin{equation}
 F_M^b(r)=f_M^b(u_r)+f_M^b(v_r).
 \label{eq:altfreq}
\end{equation}
The words are distinct, including when $r=0$.

\begin{lemma}[Alternation forces a residual excursion]
\label{lem:excursion}
Suppose $0\le a_j\le1/2$ for all $j$.
If the $r$ pairs starting at pair index $j$ are all $01$, or are all
$10$, then
\begin{equation}
 |E_{b,a}(j)|+|E_{b,a}(j+r)|\ge r/2.
 \label{eq:excursion}
\end{equation}
In particular, at least one of the two endpoint magnitudes is at
least $r/4$.
\end{lemma}
\begin{proof}
Across the block, $S_b(j+r)-S_b(j)$ is $-r$ or $r$,
while $A(j+r)-A(j)\in[0,r/2]$.
Thus $|E_{b,a}(j+r)-E_{b,a}(j)|\ge r/2$.
The triangle inequality gives \eqref{eq:excursion};
if both endpoints had magnitude less than $r/4$, it would fail.
\end{proof}

\begin{theorem}[Finite ordinary-frequency bound]
\label{thm:finite}
For every binary stream $b$, every step sequence satisfying
$0\le a_j\le1/2$, every $r\in\N$, and every $N\ge1$,
\begin{equation}
 \boxed{\displaystyle
 F_{2N}^b(r)\le
 2T_{E_{b,a}}(r/4,N)+\frac{r+1/2}{N}.}
 \label{eq:finite}
\end{equation}
\end{theorem}
\begin{proof}
Let $I_j$ be the indicator that the $r$ pairs beginning at pair index
$j$ have one of the two constant alternating orientations, and put
$K=\sum_{j<N}I_j$. Set
$q_j=\ind_{\{|E_{b,a}(j)|\ge r/4\}}$ and $B=\sum_{j<N}q_j$.
Lemma~\ref{lem:excursion} gives $I_j\le q_j+q_{j+r}$.
For every sequence $g_j\in[0,1]$,
\[
 \sum_{j<N}g_{j+r}\le\sum_{j<N}g_j+r.
\]
This follows by comparing the two shifted finite sums; at most $r$
new terms, each at most $1$, are added.
Consequently $K\le2B+r$.

An occurrence of either length-$2r+1$ alternating word at the even
start $2j$ implies $I_j=1$; an occurrence at the odd start $2j+1$
implies $I_{j+1}=1$. The words are distinct, so their combined
indicator at one starting position is at most $1$.
Summing over the $2N$ starts therefore gives
\[
 2N F_{2N}^b(r)
 \le\sum_{j<N}(I_j+I_{j+1})
 \le2K+1\le4B+2r+1.
\]
Division by $2N$ proves \eqref{eq:finite}.
The $1/(2N)$ term accounts explicitly for the final shifted start.
\end{proof}

Substitution of \eqref{eq:analyticstep} and $b=b(\pi)$ yields
\eqref{eq:mainpreview}. The source declaration is
\code{PI_ALTERNATING_FREQUENCY_TAIL_BOUND}.
No analytic or digit-identification hypothesis remains in that
specialization; the tail itself is still an unevaluated quantity.

\section{The conditional non-normality endpoint}
\label{sec:closure}

\begin{theorem}[A small residual tail excludes normality]
\label{thm:closure}
Suppose $0\le a_j\le1/2$ for every $j$. If there exist a fixed
$r\in\N$, a real $\eta$, and $N_0\in\N$ such that
\begin{equation}
 2\eta<4^{-r},\qquad
 T_{E_{b,a}}(r/4,N)\le\eta
 \quad\hbox{for all }N\ge\max(1,N_0),
 \label{eq:open}
\end{equation}
then $b$ is not binary normal.
\end{theorem}
\begin{proof}
If $b$ were normal, each of the two fixed words of length $2r+1$
would have frequency tending to $2^{-(2r+1)}$. Hence
$F_{2N}^b(r)\to4^{-r}$.
Theorem~\ref{thm:finite} and \eqref{eq:open} give, eventually,
\[
 F_{2N}^b(r)\le2\eta+\frac{r+1/2}{N}.
\]
Taking limits yields $4^{-r}\le2\eta$, a contradiction.
\end{proof}

\begin{corollary}[The actual-$\pi$ endpoint]
\label{cor:pi}
If \eqref{eq:open} holds for the residual $E_\pi$ from
\eqref{eq:analyticstep}, then $\neg\Normal(\pi)$.
\end{corollary}
\begin{proof}
The analytic step bounds and the identification of \code{piBits}
with \eqref{eq:bits} are already proved. Apply
Theorem~\ref{thm:closure}.
\end{proof}

This is the precise conditional result recorded by
\code{PI_NOT_NORMAL_OF_SMALL_ANALYTIC_RESIDUAL_TAIL}.
The eventual estimate \eqref{eq:open} has not been established for
$E_\pi$ in this repository.
For a nonvacuous search one may take $r\ge1$ and $\eta\ge0$.
The formal theorem does not need those extra assumptions:
when $r=0$, the threshold is zero and $T_E(0,N)=1$ for $N>0$,
so \eqref{eq:open} cannot hold. A negative $\eta$ is likewise
incompatible with an eventual nonnegative tail.

The quantifier order is essential. One chooses a single word length,
a single positive gap $4^{-r}-2\eta$, and then proves a bound for
every sufficiently large sample size. Letting $r$ grow with $N$,
or observing a deficit in a finite computation, does not establish
this hypothesis. The theorem is a sufficient criterion, not a claimed
characterization of all non-normal streams.

\section{What the coefficient does and does not supply}
\label{sec:limits}

\begin{proposition}[Cancellation and fixed-length analytic mass]
\label{prop:mass}
For $L\in\N$ and $M>0$,
\begin{align}
 C_t+C_{t+1}&=0,&
 \sum_{t<L}C_t&=\frac{1-(-1)^L}{8},&
 \sum_{t<L}(-1)^tC_t&=\frac L4, \label{eq:cancel}\\
 0\le\sum_{t<L}\delta(M+t)&\le\frac{L}{4M}.
 \label{eq:mass}
\end{align}
\end{proposition}
\begin{proof}
The first line is the finite geometric sum of the signed values
in \eqref{eq:Ct}; multiplying each term by $(-1)^t$ gives $1/4$.
For the second line, Lemma~\ref{lem:delta} gives
$0\le\delta(M+t)\le1/[4(M+t)]\le1/(4M)$; sum over $t<L$.
\end{proof}

Thus the aligned coefficient sum grows with $L$, but its actual
unscaled analytic correction retains a denominator. At each fixed
$L$, the right-hand side of \eqref{eq:mass} tends to zero as
$M\to\infty$.
The formal analytic word-product model makes this distinction
quantitative. For arbitrary fixed signs $\sigma_k\in\{-1,1\}$,
\begin{equation}
 \left|\prod_{k<L}\bigl(\tfrac12+\sigma_k\delta(n+k+1)\bigr)
          -2^{-L}\right|\le\frac{L}{4(n+1)}.
 \label{eq:product}
\end{equation}
Each factor lies in $[0,1]$ and differs from $1/2$ by at most
$1/[4(n+1)]$; expanding the difference of products, or inducting
on $L$, gives the bound. These analytic products and their
Ces\`aro averages converge to $2^{-L}$ at fixed $L$.
They have not been identified with actual digit frequencies.

Likewise, fixed-offset convergence in Theorem~\ref{thm:coefficient}
is not uniform in all offsets. The historical uniform-tracking
branch instead uses the physical scaling
$(N_a+t)\delta(N_a+t)$.
Neither choice of scaling constructs a finite observation of the
actual digit stream merely by defining an analytic formula.

\section{Historical routes and verification boundary}
\label{sec:verification}

The repository preserves seven historical developments alongside the
primary \code{PI_DYADIC_CT_v0_7_0} project.
Their module roots are isolated because several reuse the same Lean
module names for different definitions. Three distinctions matter
when interpreting their conditional endpoints.

First, a finite-observation route asks for an observable determined by
a fixed finite word of the actual stream at each scale and tracking
the alternating coefficient uniformly in the starting position.
Its formal frontier is equivalent to the existence of one word width
that determines start parity at every position. This digit property
remains an explicit hypothesis.

Second, the exact-obstruction branch defines a proposition
$\operatorname{Ct}(w)$ as failure of the fair-frequency limit for
the word $w$. Its equivalence
$\exists w\,\operatorname{Ct}(w)\ \Longleftrightarrow\ \neg\Normal(\pi)$
is a logical reformulation of the target.
The proposition $\operatorname{Ct}(w)$ is not the real coefficient
$C_t$, and the equivalence supplies no word witness.

Third, the legacy Part B address condition requires
$2j=2^{10+a(j)}+t(j)$ for every $j\in\N$.
At $j=0$ this is impossible. The separate audited declaration
\code{ANCHOPI.PartB.Audit.ctDigitBridge_impossible} records that
obstruction; the associated historical wrapper is vacuous.
The generic bounded-remainder argument is a different result,
and the later eventual-address formulation removes this particular
zero-index obstruction.

\subsection{A fixed, reproducible verification baseline}

This manuscript describes source commit
\begin{center}\small\url{1ce84cb63dc48e43ed92901c29a60f6d07ed4d9b}\end{center}
of \repo. Its GitHub Actions run
\href{https://github.com/randogelist/PI_NOT_NORMAL_BASE_2/actions/runs/40175464758}{40175464758}
completed successfully on 25 October 2026.
The preceding run
\href{https://github.com/randogelist/PI_NOT_NORMAL_BASE_2/actions/runs/40174579963}{40174579963}
freshly compiled all 71 modules; the baseline run reused all 71 only
after source, dependency, compiler, and output fingerprint checks,
and reran every declaration audit.
No Lean compilation for this verification baseline was delegated to
a local manuscript-building process.

\begin{center}
\begin{tabular}{lrr}
\toprule
Project root & Modules & Audited declarations\\
\midrule
Primary \code{PiCt} & 28 & 582\\
\code{branches/exact-obstruction} & 6 & 112\\
\code{branches/explicit-frontier} & 4 & 185\\
\code{branches/full-closure} & 14 & 267\\
\code{branches/part-a} & 3 & 69\\
\code{branches/part-b} & 5 & 175\\
\code{branches/reconstructed} & 5 & 256\\
\code{branches/single-hypothesis} & 6 & 113\\
\midrule
Total & 71 & 1,759\\
\bottomrule
\end{tabular}
\end{center}

Lean is pinned to version 4.33.0 and Mathlib to
\begin{center}\small\url{db584cd6d46c92f209a44c0f1c829460d327499d}.\end{center}
The hosted workflow checks the official Lean archive checksum, tests
rejection of unsupported assumptions, and audits transitive axiom
dependencies of project declarations. Its allowed axioms are
\code{propext}, \code{Classical.choice}, and \code{Quot.sound};
\code{sorryAx} and project research axioms are rejected.
The count above includes definitions and generated declarations;
it is not a count of 1,759 independent mathematical theorems.

The verification receipt and exact run identities are recorded in
\code{paper/verification-baseline.json}. The theorem map below links
the mathematical statements to the checked sources. Lean verifies
those declarations, rather than the typography or every sentence
of this exposition. In particular, compilation of an implication
does not prove its research premise.
The older Gregory/BETA9/AZETON continuation preserved in
\code{provenance/PI_CT_SOURCE_RESEARCH_v1.tex} is paper-level work
and is outside this manuscript's asserted Lean result set.

\appendix
\section{Statement-to-source map}
\label{sec:map}

Paths below are relative to the verified repository snapshot~\cite{repo}.
The complete qualified identifiers and machine-readable correspondence
are also recorded in \code{paper/theorem-map.json}.

\small
\begin{longtable}{@{}>{\raggedright\arraybackslash}p{.27\textwidth}>{\raggedright\arraybackslash}p{.68\textwidth}@{}}
\toprule
Statement & Source module and principal declaration\\
\midrule
\endfirsthead
\toprule
Statement & Source module and principal declaration\\
\midrule
\endhead
Anchor and ray readout &
\code{PiCt/Representative.lean}\newline
\code{primitive_one}, \code{primitive_ray_im},
\code{splitReadout_eq_half_add_delta}\\[5pt]
Lemma~\ref{lem:delta} &
\code{PiCt/LogCore.lean}\newline
\code{delta_bounds}, \code{scaled_bounds}\\[5pt]
Theorem~\ref{thm:coefficient} &
\code{PiCt/Instance.lean}\newline
\code{PI_CT_INSTANTIATED}, \code{PI_CT_VALUE},
\code{PI_CT_NE_ZERO}, \code{PI_CT_REMAINDER},
\code{PI_CT_CAUCHY}, \code{PI_CT_UNIQUE};
\code{PiCt/Dyadic.lean}: \code{coefficient_error}\\[5pt]
Actual digits and normality &
\code{PiCt/BinaryWords.lean}\newline
\code{piBits_eq}, \code{BinaryNormal}\\[5pt]
Equation~\eqref{eq:harmonic} &
\code{PiCt/AnalyticDrift.lean}\newline
\code{drift_error_tight}, \code{drift_error}\\[5pt]
Lemma~\ref{lem:excursion} &
\code{PiCt/AlternatingExcursions.lean}\newline
\code{residual_excursion}, \code{alternating_endpoint_tail}\\[5pt]
Theorem~\ref{thm:finite} &
\code{PiCt/AlternatingTailClosure.lean}\newline
\code{totalFreq_le_residual_tail}\\[5pt]
Actual-$\pi$ finite bound &
\code{PiCt/AlternatingTailClosure.lean}\newline
\code{PI_ALTERNATING_FREQUENCY_TAIL_BOUND}\\[5pt]
Theorem~\ref{thm:closure} and Corollary~\ref{cor:pi} &
\code{PiCt/AlternatingTailClosure.lean}\newline
\code{not_normal_of_eventually_small_residual_tail},
\code{PI_NOT_NORMAL_OF_SMALL_ANALYTIC_RESIDUAL_TAIL}\\[5pt]
Proposition~\ref{prop:mass} &
\code{PiCt/CoefficientAccumulation.lean}\newline
\code{C_adjacent_cancel}, \code{C_prefix_sum},
\code{C_aligned_prefix_sum}, \code{blockMass_bounds}\\[5pt]
Equation~\eqref{eq:product} &
\code{PiCt/AnalyticWordProduct.lean}\newline
\code{wordProduct_error}, \code{wordProduct_fixed_length_limit},
\code{wordProduct_average_limit}\\
\bottomrule
\end{longtable}
\normalsize

\begin{thebibliography}{9}
\bibitem{repo}
\textit{PI\_NOT\_NORMAL\_BASE\_2}.
Source snapshot \code{1ce84cb6}, 25 October 2026.
\href{https://github.com/randogelist/PI_NOT_NORMAL_BASE_2/tree/1ce84cb63dc48e43ed92901c29a60f6d07ed4d9b}{GitHub source snapshot}.

\bibitem{primary}
\textit{PI\_DYADIC\_CT\_v0\_7\_0}.
Recovered primary source package, 23 October 2026.
Its source identities, original hashes, and import adaptations are
preserved in the repository's \code{provenance/} directory.

\bibitem{ci}
\textit{Lean $C_t$ verification}.
GitHub Actions runs 40174579963 and 40175464758,
25 October 2026. Exact artifact digests and project counts are
retained in \code{paper/verification-baseline.json}.

\bibitem{frontier}
\textit{Proof frontier} and \textit{Source audit}.
Repository documents \code{docs/PROOF_FRONTIER.md} and
\code{docs/SOURCE_AUDIT.md}, at the source snapshot in~\cite{repo}.
\end{thebibliography}
\end{document}
